\documentclass{article}
\usepackage{iclr2027_conference,times}
%%%%% NEW MATH DEFINITIONS %%%%%

\usepackage{amsmath,amsfonts,bm}

% Mark sections of captions for referring to divisions of figures
\newcommand{\figleft}{{\em (Left)}}
\newcommand{\figcenter}{{\em (Center)}}
\newcommand{\figright}{{\em (Right)}}
\newcommand{\figtop}{{\em (Top)}}
\newcommand{\figbottom}{{\em (Bottom)}}
\newcommand{\captiona}{{\em (a)}}
\newcommand{\captionb}{{\em (b)}}
\newcommand{\captionc}{{\em (c)}}
\newcommand{\captiond}{{\em (d)}}

% Highlight a newly defined term
\newcommand{\newterm}[1]{{\bf #1}}


% Figure reference, lower-case.
\def\figref#1{figure~\ref{#1}}
% Figure reference, capital. For start of sentence
\def\Figref#1{Figure~\ref{#1}}
\def\twofigref#1#2{figures \ref{#1} and \ref{#2}}
\def\quadfigref#1#2#3#4{figures \ref{#1}, \ref{#2}, \ref{#3} and \ref{#4}}
% Section reference, lower-case.
\def\secref#1{section~\ref{#1}}
% Section reference, capital.
\def\Secref#1{Section~\ref{#1}}
% Reference to two sections.
\def\twosecrefs#1#2{sections \ref{#1} and \ref{#2}}
% Reference to three sections.
\def\secrefs#1#2#3{sections \ref{#1}, \ref{#2} and \ref{#3}}
% Reference to an equation, lower-case.
\def\eqref#1{equation~\ref{#1}}
% Reference to an equation, upper case
\def\Eqref#1{Equation~\ref{#1}}
% A raw reference to an equation---avoid using if possible
\def\plaineqref#1{\ref{#1}}
% Reference to a chapter, lower-case.
\def\chapref#1{chapter~\ref{#1}}
% Reference to an equation, upper case.
\def\Chapref#1{Chapter~\ref{#1}}
% Reference to a range of chapters
\def\rangechapref#1#2{chapters\ref{#1}--\ref{#2}}
% Reference to an algorithm, lower-case.
\def\algref#1{algorithm~\ref{#1}}
% Reference to an algorithm, upper case.
\def\Algref#1{Algorithm~\ref{#1}}
\def\twoalgref#1#2{algorithms \ref{#1} and \ref{#2}}
\def\Twoalgref#1#2{Algorithms \ref{#1} and \ref{#2}}
% Reference to a part, lower case
\def\partref#1{part~\ref{#1}}
% Reference to a part, upper case
\def\Partref#1{Part~\ref{#1}}
\def\twopartref#1#2{parts \ref{#1} and \ref{#2}}

\def\ceil#1{\lceil #1 \rceil}
\def\floor#1{\lfloor #1 \rfloor}
\def\1{\bm{1}}
\newcommand{\train}{\mathcal{D}}
\newcommand{\valid}{\mathcal{D_{\mathrm{valid}}}}
\newcommand{\test}{\mathcal{D_{\mathrm{test}}}}

\def\eps{{\epsilon}}


% Random variables
\def\reta{{\textnormal{$\eta$}}}
\def\ra{{\textnormal{a}}}
\def\rb{{\textnormal{b}}}
\def\rc{{\textnormal{c}}}
\def\rd{{\textnormal{d}}}
\def\re{{\textnormal{e}}}
\def\rf{{\textnormal{f}}}
\def\rg{{\textnormal{g}}}
\def\rh{{\textnormal{h}}}
\def\ri{{\textnormal{i}}}
\def\rj{{\textnormal{j}}}
\def\rk{{\textnormal{k}}}
\def\rl{{\textnormal{l}}}
% rm is already a command, just don't name any random variables m
\def\rn{{\textnormal{n}}}
\def\ro{{\textnormal{o}}}
\def\rp{{\textnormal{p}}}
\def\rq{{\textnormal{q}}}
\def\rr{{\textnormal{r}}}
\def\rs{{\textnormal{s}}}
\def\rt{{\textnormal{t}}}
\def\ru{{\textnormal{u}}}
\def\rv{{\textnormal{v}}}
\def\rw{{\textnormal{w}}}
\def\rx{{\textnormal{x}}}
\def\ry{{\textnormal{y}}}
\def\rz{{\textnormal{z}}}

% Random vectors
\def\rvepsilon{{\mathbf{\epsilon}}}
\def\rvtheta{{\mathbf{\theta}}}
\def\rva{{\mathbf{a}}}
\def\rvb{{\mathbf{b}}}
\def\rvc{{\mathbf{c}}}
\def\rvd{{\mathbf{d}}}
\def\rve{{\mathbf{e}}}
\def\rvf{{\mathbf{f}}}
\def\rvg{{\mathbf{g}}}
\def\rvh{{\mathbf{h}}}
\def\rvu{{\mathbf{i}}}
\def\rvj{{\mathbf{j}}}
\def\rvk{{\mathbf{k}}}
\def\rvl{{\mathbf{l}}}
\def\rvm{{\mathbf{m}}}
\def\rvn{{\mathbf{n}}}
\def\rvo{{\mathbf{o}}}
\def\rvp{{\mathbf{p}}}
\def\rvq{{\mathbf{q}}}
\def\rvr{{\mathbf{r}}}
\def\rvs{{\mathbf{s}}}
\def\rvt{{\mathbf{t}}}
\def\rvu{{\mathbf{u}}}
\def\rvv{{\mathbf{v}}}
\def\rvw{{\mathbf{w}}}
\def\rvx{{\mathbf{x}}}
\def\rvy{{\mathbf{y}}}
\def\rvz{{\mathbf{z}}}

% Elements of random vectors
\def\erva{{\textnormal{a}}}
\def\ervb{{\textnormal{b}}}
\def\ervc{{\textnormal{c}}}
\def\ervd{{\textnormal{d}}}
\def\erve{{\textnormal{e}}}
\def\ervf{{\textnormal{f}}}
\def\ervg{{\textnormal{g}}}
\def\ervh{{\textnormal{h}}}
\def\ervi{{\textnormal{i}}}
\def\ervj{{\textnormal{j}}}
\def\ervk{{\textnormal{k}}}
\def\ervl{{\textnormal{l}}}
\def\ervm{{\textnormal{m}}}
\def\ervn{{\textnormal{n}}}
\def\ervo{{\textnormal{o}}}
\def\ervp{{\textnormal{p}}}
\def\ervq{{\textnormal{q}}}
\def\ervr{{\textnormal{r}}}
\def\ervs{{\textnormal{s}}}
\def\ervt{{\textnormal{t}}}
\def\ervu{{\textnormal{u}}}
\def\ervv{{\textnormal{v}}}
\def\ervw{{\textnormal{w}}}
\def\ervx{{\textnormal{x}}}
\def\ervy{{\textnormal{y}}}
\def\ervz{{\textnormal{z}}}

% Random matrices
\def\rmA{{\mathbf{A}}}
\def\rmB{{\mathbf{B}}}
\def\rmC{{\mathbf{C}}}
\def\rmD{{\mathbf{D}}}
\def\rmE{{\mathbf{E}}}
\def\rmF{{\mathbf{F}}}
\def\rmG{{\mathbf{G}}}
\def\rmH{{\mathbf{H}}}
\def\rmI{{\mathbf{I}}}
\def\rmJ{{\mathbf{J}}}
\def\rmK{{\mathbf{K}}}
\def\rmL{{\mathbf{L}}}
\def\rmM{{\mathbf{M}}}
\def\rmN{{\mathbf{N}}}
\def\rmO{{\mathbf{O}}}
\def\rmP{{\mathbf{P}}}
\def\rmQ{{\mathbf{Q}}}
\def\rmR{{\mathbf{R}}}
\def\rmS{{\mathbf{S}}}
\def\rmT{{\mathbf{T}}}
\def\rmU{{\mathbf{U}}}
\def\rmV{{\mathbf{V}}}
\def\rmW{{\mathbf{W}}}
\def\rmX{{\mathbf{X}}}
\def\rmY{{\mathbf{Y}}}
\def\rmZ{{\mathbf{Z}}}

% Elements of random matrices
\def\ermA{{\textnormal{A}}}
\def\ermB{{\textnormal{B}}}
\def\ermC{{\textnormal{C}}}
\def\ermD{{\textnormal{D}}}
\def\ermE{{\textnormal{E}}}
\def\ermF{{\textnormal{F}}}
\def\ermG{{\textnormal{G}}}
\def\ermH{{\textnormal{H}}}
\def\ermI{{\textnormal{I}}}
\def\ermJ{{\textnormal{J}}}
\def\ermK{{\textnormal{K}}}
\def\ermL{{\textnormal{L}}}
\def\ermM{{\textnormal{M}}}
\def\ermN{{\textnormal{N}}}
\def\ermO{{\textnormal{O}}}
\def\ermP{{\textnormal{P}}}
\def\ermQ{{\textnormal{Q}}}
\def\ermR{{\textnormal{R}}}
\def\ermS{{\textnormal{S}}}
\def\ermT{{\textnormal{T}}}
\def\ermU{{\textnormal{U}}}
\def\ermV{{\textnormal{V}}}
\def\ermW{{\textnormal{W}}}
\def\ermX{{\textnormal{X}}}
\def\ermY{{\textnormal{Y}}}
\def\ermZ{{\textnormal{Z}}}

% Vectors
\def\vzero{{\bm{0}}}
\def\vone{{\bm{1}}}
\def\vmu{{\bm{\mu}}}
\def\vtheta{{\bm{\theta}}}
\def\va{{\bm{a}}}
\def\vb{{\bm{b}}}
\def\vc{{\bm{c}}}
\def\vd{{\bm{d}}}
\def\ve{{\bm{e}}}
\def\vf{{\bm{f}}}
\def\vg{{\bm{g}}}
\def\vh{{\bm{h}}}
\def\vi{{\bm{i}}}
\def\vj{{\bm{j}}}
\def\vk{{\bm{k}}}
\def\vl{{\bm{l}}}
\def\vm{{\bm{m}}}
\def\vn{{\bm{n}}}
\def\vo{{\bm{o}}}
\def\vp{{\bm{p}}}
\def\vq{{\bm{q}}}
\def\vr{{\bm{r}}}
\def\vs{{\bm{s}}}
\def\vt{{\bm{t}}}
\def\vu{{\bm{u}}}
\def\vv{{\bm{v}}}
\def\vw{{\bm{w}}}
\def\vx{{\bm{x}}}
\def\vy{{\bm{y}}}
\def\vz{{\bm{z}}}

% Elements of vectors
\def\evalpha{{\alpha}}
\def\evbeta{{\beta}}
\def\evepsilon{{\epsilon}}
\def\evlambda{{\lambda}}
\def\evomega{{\omega}}
\def\evmu{{\mu}}
\def\evpsi{{\psi}}
\def\evsigma{{\sigma}}
\def\evtheta{{\theta}}
\def\eva{{a}}
\def\evb{{b}}
\def\evc{{c}}
\def\evd{{d}}
\def\eve{{e}}
\def\evf{{f}}
\def\evg{{g}}
\def\evh{{h}}
\def\evi{{i}}
\def\evj{{j}}
\def\evk{{k}}
\def\evl{{l}}
\def\evm{{m}}
\def\evn{{n}}
\def\evo{{o}}
\def\evp{{p}}
\def\evq{{q}}
\def\evr{{r}}
\def\evs{{s}}
\def\evt{{t}}
\def\evu{{u}}
\def\evv{{v}}
\def\evw{{w}}
\def\evx{{x}}
\def\evy{{y}}
\def\evz{{z}}

% Matrix
\def\mA{{\bm{A}}}
\def\mB{{\bm{B}}}
\def\mC{{\bm{C}}}
\def\mD{{\bm{D}}}
\def\mE{{\bm{E}}}
\def\mF{{\bm{F}}}
\def\mG{{\bm{G}}}
\def\mH{{\bm{H}}}
\def\mI{{\bm{I}}}
\def\mJ{{\bm{J}}}
\def\mK{{\bm{K}}}
\def\mL{{\bm{L}}}
\def\mM{{\bm{M}}}
\def\mN{{\bm{N}}}
\def\mO{{\bm{O}}}
\def\mP{{\bm{P}}}
\def\mQ{{\bm{Q}}}
\def\mR{{\bm{R}}}
\def\mS{{\bm{S}}}
\def\mT{{\bm{T}}}
\def\mU{{\bm{U}}}
\def\mV{{\bm{V}}}
\def\mW{{\bm{W}}}
\def\mX{{\bm{X}}}
\def\mY{{\bm{Y}}}
\def\mZ{{\bm{Z}}}
\def\mBeta{{\bm{\beta}}}
\def\mPhi{{\bm{\Phi}}}
\def\mLambda{{\bm{\Lambda}}}
\def\mSigma{{\bm{\Sigma}}}

% Tensor
\DeclareMathAlphabet{\mathsfit}{\encodingdefault}{\sfdefault}{m}{sl}
\SetMathAlphabet{\mathsfit}{bold}{\encodingdefault}{\sfdefault}{bx}{n}
\newcommand{\tens}[1]{\bm{\mathsfit{#1}}}
\def\tA{{\tens{A}}}
\def\tB{{\tens{B}}}
\def\tC{{\tens{C}}}
\def\tD{{\tens{D}}}
\def\tE{{\tens{E}}}
\def\tF{{\tens{F}}}
\def\tG{{\tens{G}}}
\def\tH{{\tens{H}}}
\def\tI{{\tens{I}}}
\def\tJ{{\tens{J}}}
\def\tK{{\tens{K}}}
\def\tL{{\tens{L}}}
\def\tM{{\tens{M}}}
\def\tN{{\tens{N}}}
\def\tO{{\tens{O}}}
\def\tP{{\tens{P}}}
\def\tQ{{\tens{Q}}}
\def\tR{{\tens{R}}}
\def\tS{{\tens{S}}}
\def\tT{{\tens{T}}}
\def\tU{{\tens{U}}}
\def\tV{{\tens{V}}}
\def\tW{{\tens{W}}}
\def\tX{{\tens{X}}}
\def\tY{{\tens{Y}}}
\def\tZ{{\tens{Z}}}


% Graph
\def\gA{{\mathcal{A}}}
\def\gB{{\mathcal{B}}}
\def\gC{{\mathcal{C}}}
\def\gD{{\mathcal{D}}}
\def\gE{{\mathcal{E}}}
\def\gF{{\mathcal{F}}}
\def\gG{{\mathcal{G}}}
\def\gH{{\mathcal{H}}}
\def\gI{{\mathcal{I}}}
\def\gJ{{\mathcal{J}}}
\def\gK{{\mathcal{K}}}
\def\gL{{\mathcal{L}}}
\def\gM{{\mathcal{M}}}
\def\gN{{\mathcal{N}}}
\def\gO{{\mathcal{O}}}
\def\gP{{\mathcal{P}}}
\def\gQ{{\mathcal{Q}}}
\def\gR{{\mathcal{R}}}
\def\gS{{\mathcal{S}}}
\def\gT{{\mathcal{T}}}
\def\gU{{\mathcal{U}}}
\def\gV{{\mathcal{V}}}
\def\gW{{\mathcal{W}}}
\def\gX{{\mathcal{X}}}
\def\gY{{\mathcal{Y}}}
\def\gZ{{\mathcal{Z}}}

% Sets
\def\sA{{\mathbb{A}}}
\def\sB{{\mathbb{B}}}
\def\sC{{\mathbb{C}}}
\def\sD{{\mathbb{D}}}
% Don't use a set called E, because this would be the same as our symbol
% for expectation.
\def\sF{{\mathbb{F}}}
\def\sG{{\mathbb{G}}}
\def\sH{{\mathbb{H}}}
\def\sI{{\mathbb{I}}}
\def\sJ{{\mathbb{J}}}
\def\sK{{\mathbb{K}}}
\def\sL{{\mathbb{L}}}
\def\sM{{\mathbb{M}}}
\def\sN{{\mathbb{N}}}
\def\sO{{\mathbb{O}}}
\def\sP{{\mathbb{P}}}
\def\sQ{{\mathbb{Q}}}
\def\sR{{\mathbb{R}}}
\def\sS{{\mathbb{S}}}
\def\sT{{\mathbb{T}}}
\def\sU{{\mathbb{U}}}
\def\sV{{\mathbb{V}}}
\def\sW{{\mathbb{W}}}
\def\sX{{\mathbb{X}}}
\def\sY{{\mathbb{Y}}}
\def\sZ{{\mathbb{Z}}}

% Entries of a matrix
\def\emLambda{{\Lambda}}
\def\emA{{A}}
\def\emB{{B}}
\def\emC{{C}}
\def\emD{{D}}
\def\emE{{E}}
\def\emF{{F}}
\def\emG{{G}}
\def\emH{{H}}
\def\emI{{I}}
\def\emJ{{J}}
\def\emK{{K}}
\def\emL{{L}}
\def\emM{{M}}
\def\emN{{N}}
\def\emO{{O}}
\def\emP{{P}}
\def\emQ{{Q}}
\def\emR{{R}}
\def\emS{{S}}
\def\emT{{T}}
\def\emU{{U}}
\def\emV{{V}}
\def\emW{{W}}
\def\emX{{X}}
\def\emY{{Y}}
\def\emZ{{Z}}
\def\emSigma{{\Sigma}}

% entries of a tensor
% Same font as tensor, without \bm wrapper
\newcommand{\etens}[1]{\mathsfit{#1}}
\def\etLambda{{\etens{\Lambda}}}
\def\etA{{\etens{A}}}
\def\etB{{\etens{B}}}
\def\etC{{\etens{C}}}
\def\etD{{\etens{D}}}
\def\etE{{\etens{E}}}
\def\etF{{\etens{F}}}
\def\etG{{\etens{G}}}
\def\etH{{\etens{H}}}
\def\etI{{\etens{I}}}
\def\etJ{{\etens{J}}}
\def\etK{{\etens{K}}}
\def\etL{{\etens{L}}}
\def\etM{{\etens{M}}}
\def\etN{{\etens{N}}}
\def\etO{{\etens{O}}}
\def\etP{{\etens{P}}}
\def\etQ{{\etens{Q}}}
\def\etR{{\etens{R}}}
\def\etS{{\etens{S}}}
\def\etT{{\etens{T}}}
\def\etU{{\etens{U}}}
\def\etV{{\etens{V}}}
\def\etW{{\etens{W}}}
\def\etX{{\etens{X}}}
\def\etY{{\etens{Y}}}
\def\etZ{{\etens{Z}}}

% The true underlying data generating distribution
\newcommand{\pdata}{p_{\rm{data}}}
% The empirical distribution defined by the training set
\newcommand{\ptrain}{\hat{p}_{\rm{data}}}
\newcommand{\Ptrain}{\hat{P}_{\rm{data}}}
% The model distribution
\newcommand{\pmodel}{p_{\rm{model}}}
\newcommand{\Pmodel}{P_{\rm{model}}}
\newcommand{\ptildemodel}{\tilde{p}_{\rm{model}}}
% Stochastic autoencoder distributions
\newcommand{\pencode}{p_{\rm{encoder}}}
\newcommand{\pdecode}{p_{\rm{decoder}}}
\newcommand{\precons}{p_{\rm{reconstruct}}}

\newcommand{\laplace}{\mathrm{Laplace}} % Laplace distribution

\newcommand{\E}{\mathbb{E}}
\newcommand{\Ls}{\mathcal{L}}
\newcommand{\R}{\mathbb{R}}
\newcommand{\emp}{\tilde{p}}
\newcommand{\lr}{\alpha}
\newcommand{\reg}{\lambda}
\newcommand{\rect}{\mathrm{rectifier}}
\newcommand{\softmax}{\mathrm{softmax}}
\newcommand{\sigmoid}{\sigma}
\newcommand{\softplus}{\zeta}
\newcommand{\KL}{D_{\mathrm{KL}}}
\newcommand{\Var}{\mathrm{Var}}
\newcommand{\standarderror}{\mathrm{SE}}
\newcommand{\Cov}{\mathrm{Cov}}
% Wolfram Mathworld says $L^2$ is for function spaces and $\ell^2$ is for vectors
% But then they seem to use $L^2$ for vectors throughout the site, and so does
% wikipedia.
\newcommand{\normlzero}{L^0}
\newcommand{\normlone}{L^1}
\newcommand{\normltwo}{L^2}
\newcommand{\normlp}{L^p}
\newcommand{\normmax}{L^\infty}

\newcommand{\parents}{Pa} % See usage in notation.tex. Chosen to match Daphne's book.

\DeclareMathOperator*{\argmax}{arg\,max}
\DeclareMathOperator*{\argmin}{arg\,min}

\DeclareMathOperator{\sign}{sign}
\DeclareMathOperator{\Tr}{Tr}
\let\ab\allowbreak

\usepackage{hyperref}
\usepackage{url}
\usepackage{graphicx}
\graphicspath{{../figures/}}
\usepackage{booktabs}
\usepackage{amsmath,amssymb}
\usepackage{xcolor}
\newcommand{\todo}[1]{{\color{red}[TODO: #1]}}
\makeatletter\newcommand{\pfig}[1]{\@ifundefined{r@#1}{\todo{figure #1}}{Fig.~\ref{#1}}}\makeatother

\title{Model Width Speeds Fitting but Sets No Single Delay: A Falsification Study of Width and Delayed Generalisation on Modular Arithmetic}
\author{Anonymous authors \\ Paper under double-blind review}

\begin{document}
\maketitle

\begin{abstract}
Small networks trained on modular arithmetic often fit the training set early and generalise much later (grokking). How model size enters is unclear, because width is entangled with optimisation speed, capacity and parametrisation. We state a falsifiable account, H*, in which width gates memorisation and weight decay slowly selects a generalising circuit, so that larger models show longer delays, and test its predictions with 1\,852 runs of one-hidden-layer MLPs and one-block transformers (widths 4--256, mostly $p=59$, censoring-aware statistics). Four findings. (i) The fitting part of H* holds: the fit time falls as width$^{-0.8\ldots-1.1}$ at a fixed learning rate and standard parametrisation (a calibration of that setting), and models below $d\approx8$ cannot fit. (ii) Larger models generalise from \emph{less} data: the smallest train fraction at which the transformer generalises falls from 0.5 to 0.35 over widths 16 to 128, and for the MLP from 0.35 to 0.25 over widths 16 to 256 with a 100k-step cap (6 seeds per cell; at shorter caps the trend is hidden by censoring). (iii) The generalisation time follows no single law in width (nearly flat for MLP addition at train fraction 0.5, decreasing at lower fractions and for $a^2+b^2$, inconclusive for the transformer once two accuracy thresholds must agree), and the delay is non-monotone or rises only because the fit time falls, so the part of H* that ties the delay to capacity is not supported. (iv) The one clean width trend of the generalisation time we isolated (weight decay 3, fixed learning rate: log--log slope $-0.40$, 95\% CI $[-0.53,-0.27]$) is not independent of optimisation speed: raising only the learning rate per width (to 30--100$\times$ the baseline) removes 97--99\% of the delay, and the width trend then goes undetected at train fractions 0.4 and 0.5 but survives at 0.3 (carried by the smallest width), so the effect of tuning is inconclusive and depends on the train fraction. A post hoc conjecture that width sets only the fit time is falsified on held-out settings. Scope: width (plus a depth-2 control), mainly $p=59$ with addition, finite step caps.
\end{abstract}

\section{Introduction}
Neural networks trained on small algorithmic tasks such as modular addition often reach near-perfect training accuracy early, while validation accuracy rises only much later---\emph{delayed generalization}, or grokking \citep{power2022grokking}. Explanations so far emphasise circuit efficiency under weight decay \citep{varma2023explaining}, a transition from lazy to rich training dynamics \citep{kumar2024grokking}, weight-norm landscapes \citep{liu2023omnigrok}, and the gradual formation of Fourier-feature circuits followed by clean-up of memorising components \citep{nanda2023progress}. How the \emph{size} of the model enters is less systematically characterised: size is entangled with parametrisation, initialisation scale, optimisation speed and the capacity to memorise the training set.

We state a falsifiable account (H*, Sec.~\ref{sec:method}), derive six predictions that can fail (P1, P1b, P2--P6), and test them with multi-seed experiments with censoring-aware statistics. Our contributions are: (i) a set of predictions written in a versioned study plan, most of them before the main sweeps (the history and the deviations D1--D5 are in Sec.~\ref{sec:method}) in which fit time, generalisation time and their difference are analysed separately, because the delay $\Delta=t_{\mathrm{gen}}-t_{\mathrm{fit}}$ can grow with width only because $t_{\mathrm{fit}}$ shrinks; (ii) a width sweep over two architectures with censoring kept in every statistic, plus a phase diagram over width and train fraction, a random-label capacity calibration, a learning-rate/weight-decay time-scale test, a learning-rate-only tuning test at fixed weight decay, and scope checks over operation, prime and depth; (iii) a verdict for each prediction, and for a seventh, post hoc conjecture (P7) that the first results suggested and that held-out settings falsified; (iv) a decision rule, declared after the first review, that a verdict about a width trend counts only if its class agrees at the 0.99 and 0.95 accuracy thresholds. H* as a whole would have been confirmed by a fit time that falls with width, a gate near the capacity needed to fit random labels, a delay that does not shrink above the gate and follows a smooth law in width, and a width effect that survives tuning of the learning rate and weight decay. We find the first, the second only partly, and not the others: the part of H* that ties the \emph{length of the delay} to capacity is not supported for these models. In short, width sets the fit time robustly, but we find no single width--delay law, and the one clean width trend of the delay we could isolate (weight decay 3) is not independent of optimisation speed: a larger learning rate removes most of the delay and, at train fractions 0.4 and 0.5, the detectable trend. The finding we would defend as new is the second: the smallest train fraction at which a model generalises falls with width, and the MLP result appears only with a step cap long enough (Sec.~4.3, \pfig{fig:fgen}); the fit-time law is a calibration of the setting.

\section{Related Work}
\textbf{Grokking on algorithmic data.} \citet{power2022grokking} introduced binary-operation datasets modulo a prime and reported delayed generalization, with weight decay the most effective intervention for data efficiency. \citet{nanda2023progress} reverse-engineered a one-layer transformer on modular addition and split training into memorisation, circuit formation and clean-up, with restricted/excluded loss and Fourier-component measures as progress measures.
\textbf{Explanations.} \citet{varma2023explaining} attribute grokking to a slower-to-learn but more efficient generalising circuit and predict a critical data size; \citet{kumar2024grokking} tie it to lazy-to-rich dynamics, with an output-scale parameter controlling it; \citet{liu2023omnigrok} relate it to weight norm and initialisation scale, showing initialisation scale is a confound for any size study.
\citet{liu2022towards} give an effective theory and phase diagrams over hyperparameters and training-set size with a ``Goldilocks zone'' between memorisation and confusion; \citet{gromov2023grokking} gives analytic solutions for two-layer networks on modular arithmetic; \citet{thilak2022slingshot} tie grokking to late-training optimiser instabilities visible in the last-layer norm; \citet{zhong2023clock} show that small changes to hyperparameters and initialisation produce different algorithms on the same modular task; \citet{merrill2023tale} describe the transition as the emergence of a sparse subnetwork in sparse parity. These works vary data size, optimiser or initialisation; none, to our reading of their abstracts, isolates model width against fit and generalisation times, which is the target here.
\textbf{Capacity, memorisation and double descent.} \citet{zhang2017understanding} show networks fit random labels; \citet{nakkiran2020deep} report model-wise and epoch-wise double descent; \citet{davies2023unifying} relate grokking to double descent through the relative learning speed of patterns; \citet{morris2025memorize} estimate memorisation capacity per parameter in language models and relate the onset of generalisation to capacity filling; \citet{tirumala2022memorization} study memorisation dynamics in language models. We use AdamW \citep{loshchilov2017decoupled} throughout.
\textbf{Positioning.} \citet{davies2023unifying} already report model-wise grokking for modular division at $p=97$ with a two-layer transformer and a small weight decay ($10^{-5}$), so the existence of a size dependence is not new, and a fit time that falls with width at fixed learning rate and standard parametrisation is plausible under standard parametrisation (wider layers receive larger effective updates at a fixed learning rate), so we treat it as a calibration of the setting rather than a discovery and do not claim it for other parametrisations. What is new is the measurement and the negative results: we quantify fit time, generalisation time and their difference separately with seed-level uncertainty and censoring in the weight-decay-1 regime of \citet{power2022grokking} and \citet{nanda2023progress}; we cross width with train fraction, operation, prime and architecture; we control the learning-rate and weight-decay time scales with per-width tuning on held-out seeds; and we report which parts of a capacity-based account survive. The questions of mechanism and initialisation scale that the weight-norm and lazy-to-rich accounts raise are probed only by small checks (Sec.~4.5).

\section{Method}\label{sec:method}
\paragraph{Data.} For a prime $p$ (primary: $p=59$; $p=97$ for scope checks) the task is $a+b \bmod p$ over all $p^2$ ordered pairs, with a random fraction $f$ as the training set and the remainder as validation. The seed controls both the initialisation and the split; further operations (subtraction, $a^2+b^2$) are used in the scope experiments. A fixed random labelling of all $p^2$ inputs, with the same split fractions, is used to measure memorisation capacity (E0b).

\paragraph{Models and size.} \texttt{tfm}: a one-layer pre-residual transformer without LayerNorm, 4 heads, MLP width $4d$, learned positions, $\mathcal N(0,1/d)$ embedding initialisation, read out at the ``='' token. \texttt{mlp}: concatenated embeddings of $a$ and $b$, one ReLU hidden layer of width $4d$. Size is the model width $d$; we also report the parameter counts with and without embeddings, and parameters relative to $fp^2$, because width, parameter count and capacity relative to the training set are different notions of size.

\paragraph{Optimisation.} Full-batch AdamW \citep{loshchilov2017decoupled} ($\beta=(0.9,0.98)$), learning rate $10^{-3}$, weight decay $1$ on all parameters unless swept, in line with the settings of \citet{power2022grokking} and \citet{nanda2023progress}. Step budgets are 25\,000 (\texttt{tfm}) and 20\,000 (\texttt{mlp}); a run stops 500 steps after validation accuracy first reaches 0.99. Metrics are logged every 100 (\texttt{tfm}) or 50 (\texttt{mlp}) steps (5 in E9, 25 in E7, 250 in E8).

\paragraph{Operational definitions.} $t_{\mathrm{fit}}$ and $t_{\mathrm{gen}}$ are the first logged steps from which train, respectively validation, accuracy stays $\ge 0.99$ for three consecutive evaluations. The delay is $\Delta=t_{\mathrm{gen}}-t_{\mathrm{fit}}$, defined when the run fits the training set. A run that fits but never generalises within the budget is right-censored at the steps run and is kept, never dropped. Runs are classed as no fit, memorise-only, grokking ($\Delta>5$ evaluation intervals) or direct generalisation. Sensitivity analyses use the 0.95 threshold and the 50\% validation crossing.

\paragraph{Hypothesis H* and predictions.} H*: delayed generalisation needs a memorising solution that a model of sufficient capacity reaches quickly and a generalising (Fourier) circuit that weight decay selects slowly, in the spirit of the circuit-efficiency and lazy-to-rich accounts \citep{varma2023explaining,kumar2024grokking} and of the capacity-filling account \citep{morris2025memorize}. Size then acts through whether memorisation is possible and how fast it happens. The predictions, each of which can fail, are:
\textbf{P1} (fit gate): below a width $w_c(f)$ models do not fit the training set, $w_c$ grows with the training-set size, and delays do not occur far below the capacity needed to fit random labels of the same set. \textbf{P1b}: the minimum train fraction $f_{\mathrm{gen}}(w)$ at which runs generalise falls with width (falsified if flat over $\ge5$ widths). \textbf{P2}: $t_{\mathrm{fit}}$ falls with width, and $\Delta$ is non-decreasing above the gate; H* is falsified by decreasing $\Delta$, by a decreasing or flat $t_{\mathrm{gen}}$ with $\Delta$ rising only through $t_{\mathrm{fit}}$, or by flat $\Delta$ over a decade of width (flat means a log--log slope CI inside $\pm0.15$). \textbf{P3}: if $\Delta$ depends on width the trend is a smooth monotone power law. \textbf{P4}: (i) steps$\cdot$lr$\cdot$wd is the natural clock, and (ii) a width trend survives per-width tuning of lr and wd; if it vanishes, size acts only through optimisation speed. \textbf{P5}: near the lowest generalising $f$ larger models show longer or lost grokking, at large $f$ delays shrink for all widths. \textbf{P6}: the patterns carry over to a second operation, architecture and prime. \textbf{P7} (post hoc, registered after the MLP width sweep at $f=0.5$ and tested only on settings not used to form it): width sets $t_{\mathrm{fit}}\propto w^{-1}$ but not $t_{\mathrm{gen}}$, whose log--log slope stays inside $\pm0.15$, also after tuning.
Allowed verdicts are supported, refuted or conditional; none was fixed in advance. The rival accounts we try to discriminate are circuit efficiency, capacity filling, lazy-to-rich dynamics, a weight-norm ``Goldilocks'' zone \citep{liu2023omnigrok}, double descent \citep{davies2023unifying,nakkiran2020deep} and the relative learning speed of memorising and generalising features.

\paragraph{Statistics.} Seeds are the unit of replication (seed fixes initialisation and split): 16 per width for the transformer sweep ($8$ at $d=256$), 8 for the MLP, 3 in the phase diagram (5--8 near the transition), 3 plus 3 fresh seeds in the tuning experiments, 3--10 in the scope checks. We report medians with 95\% bootstrap CIs over seeds, Kaplan--Meier medians for censored data, and the slope $\beta$ of $\log t$ on $\log d$ from a lognormal Tobit regression, which keeps censored runs, with 95\% bootstrap CIs over runs (resampling whole seeds gives the same class in 40 of 43 setting--threshold pairs); Spearman correlation between width and $\Delta$ uses a permutation test. A slope is \emph{nonzero} if its CI excludes 0, \emph{flat} if the CI lies inside $\pm0.15$ (a change of at most 29\% per decade of width, a convention fixed before the sweeps) and \emph{inconclusive} otherwise; for each setting we also report the minimum detectable $|\beta|$ at 80\% power (2.8 design standard errors), which is $0.1$--$0.4$ for the MLP cells and $0.2$--$0.7$ for 3-seed transformer cells, so ``inconclusive'' for a transformer means no power rather than no effect. After the first review we adopted a decision rule: a verdict about a width trend counts only if its class is the same at the 0.99 and the 0.95 threshold (the 50\% crossing is reported as a third check); 19 of the 24 settings of the main table meet it (Table~\ref{tab:threshold_rule}, Appendix).

The plan is delivered as \texttt{plan/PLAN.md} with its version history. P1 and P2--P6 with their falsifiers were committed before the first training code and any result; P1b, the separate treatment of $t_{\mathrm{fit}}$, $t_{\mathrm{gen}}$ and $\Delta$ in P2, and the 0.95 co-primary threshold were added after a two-seed pilot, once the first sweeps were running, so they are not pre-stated in the strict sense; P7, the scope settings, E5-lite, E6, the decision rule and the follow-ups E7--E13 are post hoc or exploratory (exact times and deviations D1--D5: Appendix~\ref{app:plan}). No multiplicity correction is applied; a lone interval that barely excludes zero or the $\pm0.15$ band is treated as suggestive, and claims rest on several concordant settings.

\section{Experiments}
The experiments are labelled as in the code and results folders: E0b random-label capacity (P1), E1 width sweeps at $f=0.5$ (P2, P3), E2 and E2b the width $\times$ train-fraction phase diagram (P1b, P5), E3 the learning-rate/weight-decay grid and tuning test (P4), E4 scope checks (P6), E5-lite the weight-norm analysis of logged curves, E6 the MLP initialisation-scale sweep, and, added after the first review, E7 learning-rate-only tuning at fixed weight decay, E8 and E8b a long-cap, cap-matched low-$f$ grid, E9 fine-resolution $t_{\mathrm{fit}}$, E10 extra seeds, E11 a reproducibility check, E12 and E12b Fourier structure, E13 a depth-2 control and, after a second review, E14 lr tuning away from the logging floor. $\Delta=t_{\mathrm{gen}}-t_{\mathrm{fit}}$, $f$ is the train fraction and $d$ the width. All experiments run on CPU; $p=59$, $f=0.5$, addition, MLP or transformer, AdamW with lr $10^{-3}$ and weight decay 1 unless stated. Every number below is read from a file in \texttt{results/} (per-line sources are in \texttt{results/verdicts.md} and \texttt{results/followups.md}). A ``step cap'' bounds every run (20--125k steps), so ``never generalises'' always means ``not within the cap''.

\subsection{Fit time scales as 1/width; generalisation time does not follow one law (P2, P3)}
Across the E1 width sweep, $t_{\mathrm{fit}}$ falls with width with log--log slope $-0.97$ (95\% CI $[-1.04,-0.90]$) for the MLP and $-0.83$ $[-0.87,-0.78]$ for the transformer; across the phase diagram the slopes lie in $-0.85\ldots-1.01$ (MLP) and $-0.71\ldots-1.2$ (transformer), in the learning-rate experiment at $-1.12$ and $-0.92$, and with two hidden layers (E13, 3 seeds) at $-0.97$ $[-1.10,-0.83]$ and $-1.22$ $[-2.18,-0.40]$. The same holds for subtraction, $a^2+b^2$ and $p=97$, so P2 is supported for $t_{\mathrm{fit}}$. Logging every 5 steps (E9, 5 seeds, $d=16$--$256$) gives slopes $-1.06$ $[-1.10,-1.03]$ (MLP) and $-0.91$ $[-0.95,-0.85]$ (transformer), and the coarse logging of E1 overstated $t_{\mathrm{fit}}$ at $d\ge128$ by 18--43\% (MLP medians 1420, 545, 240, 125, 75 steps). The law is only approximate: the local exponent is $-1.29$ $[-1.36,-1.21]$ (MLP) and $-1.03$ $[-1.14,-0.87]$ (transformer) over $d=16$--$64$ and $-0.84$ and $-0.79$ over $d=64$--$256$, a partial flattening that a floor of order 50--100 steps would produce.

For $t_{\mathrm{gen}}$ the picture is not a power law (\pfig{fig:delay}). The MLP slope is $-0.10$ $[-0.22,-0.01]$ at both thresholds, with Kaplan--Meier medians of 3100, 5100, 4500, 3300 and 2800 steps from $d=16$ to $256$: non-monotone, a decrease of a factor 0.76 over 1.2 decades, close to but not inside the flat band. For the transformer the slope is $-0.28$ $[-0.57,-0.04]$ at the 0.99 threshold but $-0.06$ $[-0.32,0.18]$ at 0.95 (4 of 72 runs sit on a plateau just under 0.99, D3), so it fails the two-threshold rule and stays inconclusive (pre-specified seeds 0--7: $-0.45$ $[-0.82,-0.12]$ at 0.99, $-0.22$ $[-0.55,0.09]$ at 0.95). The delay $\Delta$ of the MLP rises from $d=16$ to $32$ (medians 1500 and 5225 steps) and then falls to 2700 (Spearman $0.09$, $p=0.59$): neither non-decreasing, as H* requires, nor decreasing. For the transformer at 0.95 $\Delta$ increases with width (Spearman $0.36$, $p=0.003$, $n=69$; medians 150 to 3150 to 1650), the pattern ``$\Delta$ rises only through $t_{\mathrm{fit}}$'' that we named as a falsifier of H*. P3 (a smooth scale law for $\Delta$) is not supported. Parameters per training example do not collapse the two families (cells in one factor-2 ratio bin range from 0 to 100\% of seeds generalising; 59 cells), because the ratio mixes the roles of $f$ and width.

\begin{figure}[t]
\centering
\includegraphics[width=\linewidth]{f2_delay_vs_size.pdf}
\caption{Fit and generalisation time against width at $f=0.5$, addition mod 59. Left: Kaplan--Meier median steps to fit (train accuracy $\ge0.99$) and to generalise (validation $\ge0.99$ and $\ge0.95$) with 95\% bootstrap CIs over seeds (MLP 8 seeds; transformer 16, and 8 at $d=256$), with log--log slopes and CIs. Right: delay $t_{\mathrm{gen}}-t_{\mathrm{fit}}$ per seed (0.99); open triangles are runs that did not generalise within the step cap (lower bounds); the line is the Kaplan--Meier median.}
\label{fig:delay}
\end{figure}

\subsection{A fit gate exists, but it does not locate the delay (P1)}
At $f=0.5$ and weight decay 1, neither family fits at $d=4$ (0 of 8 seeds). At $d=8$ the MLP fits in 6 of 8 runs but generalises in 2 of 8, with $t_{\mathrm{gen}}\approx t_{\mathrm{fit}}$, so there is no delay at the gate; the transformer fits in 4 of 8 and never reaches 0.99 validation accuracy (3 of 8 at 0.95). The smallest width that fits a fixed random labelling (\pfig{fig:capacity}) is non-decreasing in $f$ in 4 of 4 series (family $\times$ weight decay), on a width grid of factor 2. Against P1, the transformer with $d=16$ (5.1k parameters, $3\times$ below its random-label capacity) still groks: 13 of 16 runs generalise with a median delay of 2000 steps. So the random-label capacity is an upper bound on the fit gate, not the place where delays begin; P1 is only partly supported.

\subsection{Larger models generalise from less data (P1b, P5)}
In the phase diagram over width and train fraction (\pfig{fig:map}, \pfig{fig:fgen}), the minimum fraction at which at least 2/3 of the seeds generalise (also counting the 16-seed $f=0.5$ runs) is $f_{\mathrm{gen}}=0.5,0.5,0.4,0.35$ at 0.99 for transformer widths $16,32,64,128$, supporting P1b for the transformer. For the MLP, $f_{\mathrm{gen}}=0.4$ at all five widths with a 20k cap, which would falsify P1b by our criterion, but this is a cap effect. In the long-cap grid (E8/E8b: 100k steps for every cell, 6 seeds, Wilson 95\% intervals; \pfig{fig:fgen}) the fraction of seeds that generalise at $f=0.3$ is 0/6, 1/6, 6/6, 6/6, 6/6 for $d=16,32,64,128,256$ (a 0/6 cell has upper bound 0.39), at $f=0.35$ 5/6 and then 6/6, at $f=0.25$ 0/6 for $d\le64$ and 6/6 for $d=128,256$, and at $f=0.2$ 0/6 at every width. Taking $f_{\mathrm{gen}}$ as the smallest tested $f$ with at least 3 of 6 seeds generalising, $f_{\mathrm{gen}}=0.35,0.35,0.30,0.25,0.25$: it falls with width, within a 0.05 grid and a finite cap; the $f=0.25$ points need up to 76k steps, so the zeros at small $d$ and at $f=0.2$ may be a slower rather than an absent transition. The transformer shows the same direction at $f=0.3$ (100k steps): 2/6 seeds at $d=64$ and 5/6 at $d=128$, and 0/3 at $d=32$ within 125k (an earlier 3-seed reading of 1/3 and 2/3 was anecdotal). The width effect on time interacts with $f$ in the opposite direction to P5, which predicted longer or lost grokking for larger models near the lowest generalising $f$: the MLP slope of $t_{\mathrm{gen}}$ is $-0.53$ $[-0.62,-0.42]$ at $f=0.3$ ($d=64$--$256$, all cells with 6/6 events, same cap), $-0.72$ $[-0.84,-0.60]$ at $0.35$ ($d=32$--$256$), $-0.32$ $[-0.49,-0.17]$ at $0.4$ (8 seeds), $-0.23$ $[-0.45,-0.08]$ at $0.7$, and $-0.10$ at $0.5$; the transformer slopes are negative at every $f\ge0.4$ but 3-seed cells give wide intervals (minimum detectable slope 0.4--0.7). At $f=0.6$ the MLP slope is $-0.12$ $[-0.18,-0.06]$ at 0.99 and $-0.10$ $[-0.15,-0.04]$ at 0.95, inside the flat band at 0.95 only, so it fails the rule and is reported as a small decrease. Larger models were the ones that still generalised at the lowest $f$, and faster; we saw no ``lost grokking'' or double-descent-like deterioration up to $d=256$.

\begin{figure}[t]
\centering
\includegraphics[width=\linewidth]{f9_fgen_width.pdf}
\caption{Share of seeds that generalise (validation accuracy $\ge0.99$, sustained) within the step cap against train fraction, one curve per width; bars are Wilson 95\% intervals. (a) MLP, cap 100k steps, 6 seeds per cell: $f_{\mathrm{gen}}$ (smallest $f$ with at least half of the seeds generalising) falls from 0.35 at $d=16$ to 0.25 at $d=256$. (b) Transformer, cap 25k steps, 3--8 seeds per cell; crosses: $f=0.3$ at cap 100k, 6 seeds.}
\label{fig:fgen}
\end{figure}

\subsection{Controls: learning-rate and weight-decay clocks, and tuning (P4)}
Over lr $\in\{3\!\times\!10^{-4},10^{-3},3\!\times\!10^{-3}\}$ and wd $\in\{0.3,1,3\}$ (189 runs, budgets scaled by $1/(\mathrm{lr}\cdot\mathrm{wd})$), the MLP generalisation time scales as $\mathrm{lr}^{-1.05\ldots-1.10}$ and $\mathrm{wd}^{-0.63\ldots-0.93}$ (Fig.~\ref{fig:controls}, App.); rescaling by lr$\cdot$wd reduces the spread of $\log_{10}t_{\mathrm{gen}}$ over the nine cells from 0.55--0.58 to 0.07--0.19 for the MLP, but only to 0.25--0.38 for the transformer, where lr does not act as $1/\mathrm{lr}$. The lr$\cdot$wd clock is approximate for the MLP, not a product law for the transformer, and not a common clock for fitting and generalising ($t_{\mathrm{fit}}$ hardly depends on wd).

\emph{A confounded tuning test.} Picking for each width the (lr, wd) cell with the shortest median $t_{\mathrm{gen}}$ on seeds 0--2 and re-estimating on fresh seeds 3--5 gives an MLP slope of $-0.28$ $[-0.41,-0.16]$ (baseline cell: $0.06$ $[-0.26,0.29]$). It does not isolate width: the selected cells are $d=16$ at wd 1 and $d=32,64,128$ at wd 3, because $d=16$ never fits at wd 3 (0 of 27 runs, lr $3\!\times\!10^{-4}$ to 0.1). A comparison at fixed (lr, wd) removes the confound: at wd 3 and lr $3\!\times\!10^{-3}$ over $d=32,64,128$ the fresh-seed slope is $-0.40$ $[-0.53,-0.27]$ ($-0.34$ $[-0.49,-0.17]$ at 0.95), while every wd $\le1$ column is flat or inconclusive ($-0.07$ to $+0.08$). The steeper $-0.66$ to $-1.79$ of the wd 3 grid columns includes $d=16$ censored at the cap and is not a width trend.

\emph{Tuning the learning rate alone (E7, 321 runs).} At each fixed wd $\in\{1,2,3\}$ and each width we chose the lr with the smallest median $t_{\mathrm{gen}}$ on seeds 0--2 from $\{3\!\times\!10^{-4},\ldots,0.1\}$ and measured on seeds 3--5 (\pfig{fig:lrtune}a; its legend gives the $d=32$--$128$ slopes). The optimum lies at the top of the grid, lr 0.03--0.1 (the grid was extended twice because earlier optima sat on its edge), and the median $t_{\mathrm{gen}}$ is 75--150 steps at every width that fits on the fresh seeds (100--150 on the selection seeds), with a gap $t_{\mathrm{gen}}-t_{\mathrm{fit}}$ of 0--100 steps (logging step 25), against 1850--5000 steps at the baseline lr: for this MLP at $p=59$ the delay is removed by a faster learning rate. The fresh-seed slopes are $+0.15$ $[-0.23,0.53]$ (wd 1, $d=16$--$128$), $+0.05$ $[-0.14,0.23]$ (wd 2) and $-0.11$ $[-0.63,0.35]$ (wd 3, $d=32$--$128$): all inconclusive, with minimum detectable slopes of 0.26--0.60. The trend $-0.40$ at wd 3 is therefore a property of a fixed, slow learning rate and not independent of optimisation speed. This test is weak: $t_{\mathrm{gen}}$ cannot go below $t_{\mathrm{fit}}$ and the logging resolution, the optimum sits at the top of the grid, $d=128$ is slower at lr 0.1 than at 0.03 (an instability), and the transformer test (three widths, one lr, minimum detectable slope 1.2--1.7) is uninformative.

\emph{Away from the floor (E14, 336 runs).} At wd 1 and $f=0.4$ and $0.3$, where the baseline delay is 7\,000--100\,000+ steps, we extended the lr grid to $10^{-3}$--1 with 10-step logging, chose the lr per width on seeds 0--2 and measured on seeds 3--5 (\pfig{fig:lrtune}b; threshold-rule table~\ref{tab:threshold_rule_followups}, Appendix). The optimum is interior in all eight cells (lr 0.1, or 0.3 at one width; lr 1 never reaches 0.99), and a 30--100$\times$ larger lr removes 97--99\% of the delay (median $t_{\mathrm{gen}}$ at $f=0.4$: 7\,200--19\,200 steps at lr $10^{-3}$, 180--350 tuned). At $f=0.4$ the tuned slope is $-0.03$ $[-0.23,0.17]$ (0.95: $+0.01$ $[-0.18,0.22]$; minimum detectable 0.29), inconclusive rather than flat, and excludes the untuned $-0.39$ $[-0.65,-0.17]$. At $f=0.3$ a trend survives, $-0.78$ $[-1.67,-0.23]$ (0.95: $-0.47$ $[-0.89,-0.17]$), but it is carried by $d=16$ ($d=32$--$128$: $-0.30$ $[-0.65,0.02]$, inconclusive). So P4(ii) is \emph{inconclusive and depends on $f$}: no width trend is detected once the lr is tuned at $f=0.4$ and $0.5$ (the latter a weak test), but one survives at $f=0.3$. Whether size acts \emph{only} through optimisation speed is limited by one family, one task, a fixed weight decay in E14 and the transformer being untested.

\begin{figure}[t]
\centering
\includegraphics[width=0.88\linewidth]{f8_followups.pdf}
\caption{(a) MLP at $f=0.5$, evaluated every 25 steps: $t_{\mathrm{gen}}$ (0.99) against width on fresh seeds 3--5, lr $3\!\times\!10^{-3}$ fixed at weight decay 3 against the lr tuned per width (chosen on seeds 0--2) at wd 3, 2 and 1; legend: $d=32$--$128$ slope with 95\% CI; open triangles: no $t_{\mathrm{gen}}$ within the cap (the fixed-lr $d=16$ point is from seeds 0--2). (b) MLP at $f=0.4$ and $0.3$, wd 1, evaluated every 10 steps (E14): lr tuned per width on seeds 0--2 (interior of the grid) against the baseline lr $10^{-3}$; dots are seeds, lines Kaplan--Meier medians; legend: Tobit slope over $d=16$--$128$ with 95\% CI.}
\label{fig:lrtune}
\end{figure}

\subsection{Scope: fit time carries over, generalisation time does not (P6)}
The relation $t_{\mathrm{fit}}\propto1/w$ holds for both families, three operations, two primes and two depths (Fig.~\ref{fig:scope}, Appendix). The width dependence of $t_{\mathrm{gen}}$ does not: for MLP subtraction (5 seeds) the slope is $+0.07$ $[-0.21,0.32]$ (inconclusive), for $a^2+b^2$ it is $-0.78$ $[-0.95,-0.61]$ (0.95: $-0.53$; $t_{\mathrm{gen}}$ falls five-fold from $d=16$ to $128$), and for MLP addition at $p=97$ (8 seeds) it is $-0.13$ $[-0.34,0.03]$ (inconclusive; 10 seeds for subtraction in E10 give $+0.06$ $[-0.10,0.21]$, still inconclusive). Transformers at $p=97$ (3 seeds) all reach 0.99 with delays of 0--700 steps, much shorter than at $p=59$; 3 seeds are too few for a slope. With two hidden layers (E13, 3 seeds) the MLP $t_{\mathrm{gen}}$ slope is $-0.53$ $[-0.98,0.01]$ (depth 1 on the same seeds: $+0.06$ $[-0.26,0.33]$), inconclusive; size here is width at depth 1 plus this control, not parameter count.

\input{../tables/t3_verdicts.tex}

\subsection{Mechanism and initialisation checks (rivals: circuit efficiency, lazy-to-rich)}
\emph{Fourier structure (E12, E12b; 5 seeds, $d=16,64,256$, both families).} We tracked two selection-free measures: the participation ratio (PR) of the embedding's Fourier power over $\mathbb Z_{59}$ and of the logits' two-dimensional Fourier power (random: 28 and about 3\,300; a five-frequency solution: 5 and 10), and report the position $(t_{\mathrm{event}}-t_{\mathrm{fit}})/(t_{\mathrm{gen}}-t_{\mathrm{fit}})$ of the half-way point of $\log$ PR (medians; descriptive, since each run's own end points define it; Table~\ref{tab:fourier}, Appendix). For $d\ge64$ the logit spectrum is compact at or soon after the fit (position 0.00--0.14, at most 0.42 in any run); the embedding concentrates later, at 0.52--0.61 of the delay for the transformer and 0.66--0.78 for the MLP, and the wide MLP never does (PR 25 of 29 at $d=256$). At $d=16$ both are at or before the fit. This is compatible with a sparse output circuit in place soon after the fit while generalisation is still most of the delay away; it cannot show causality (no intervention; a sparse spectrum is also what a nearly generalising network looks like) or that the compact component is the one that generalises. Five key frequencies chosen from each run's last snapshot (E12) give the same ordering but are partly circular, and reproduce the full loss for the transformer and the narrow MLP only.

\emph{Initialisation scale (E6, exploratory, MLP).} Multiplying all parameters at initialisation by 0.5, 1 or 2 (5 seeds per cell, 60 runs; scale 1 reproduces E1 exactly) leaves $t_{\mathrm{fit}}$ slopes of $-1.05$, $-1.10$, $-1.20$ and inconclusive $t_{\mathrm{gen}}$ slopes ($-0.25$ to $+0.06$, CIs wider than the band). We read this as no detectable effect of a 4$\times$ range of initial scale, not a refutation of lazy-to-rich accounts.

\subsection{The post hoc conjecture P7 is falsified}
The MLP sweep at $f=0.5$ suggested that width sets $t_{\mathrm{fit}}$ only. We stated this as P7 after seeing that sweep and tested it on held-out settings. It fails where the two-threshold rule is met: the $t_{\mathrm{gen}}$ slope has a CI excluding the $\pm0.15$ band (the class is \emph{nonzero} at both thresholds) for the MLP at $f=0.3$--$0.4$ and $0.7$, for $a^2+b^2$, and for the MLP at wd 3 with a fixed learning rate. It is not falsified by the transformer at $f=0.5$ or by $f=0.6$ (both fail the rule), nor at the learning-rate optimum at $f=0.4$ and $0.5$ (E14, E7), where no trend is detected (a trend survives at $f=0.3$). P7 is falsified as a general statement; it survives only as a description of untuned MLP addition and subtraction at $f=0.5$ ($p=97$) and of the MLP at its learning-rate optimum at $f\ge0.4$, both inconclusive rather than confirmed flat.


\section{Discussion and Limitations}
\textbf{What the evidence says about the account.} The part of H* that concerns fitting is supported: models below $d\approx8$ cannot fit, and the fit time falls roughly as $1/\text{width}$ in every setting we tried, all with standard parametrisation (wider layers take larger effective steps at a fixed learning rate, so this is a calibration, not a discovery) and learning rate $10^{-3}$, with a local exponent that flattens above $d=64$. The part that links capacity to the delay is not. The random-label capacity does not locate the delay, $\Delta$ is non-monotone or rises only through $t_{\mathrm{fit}}$, and larger models generalise from less data and faster rather than later. The one clean width trend of the generalisation time that we isolated (weight decay 3, fixed lr, slope $-0.40$) goes undetected at $f=0.4$ and $0.5$ when the learning rate alone is tuned per width (97--99\% of the delay is removed) but survives at $f=0.3$, so for this MLP the width effect on the delay is not separable from optimisation speed in general, and the test result depends on $f$. Our verdict on the capacity-gate account of grokking vs.\ size is therefore \emph{conditional}: it holds as a gate on fitting, and fails as an explanation of the length of the delay. Capacity filling \citep{morris2025memorize} is consistent with the fit gate and $f_{\mathrm{gen}}(w)$. A relative-learning-speed view \citep{davies2023unifying} agrees with the MLP clocks ($t_{\mathrm{gen}}\propto \mathrm{lr}^{-1.1}\mathrm{wd}^{-0.7}$) and with E7 and E14, without being confirmed. A weight-norm clean-up view \citep{liu2023omnigrok,nanda2023progress} is not contradicted by the logged parameter norm: it rises, peaks and falls, and generalisation at 0.99 occurs once it has fallen to 0.71--0.80 of its peak (MLP $d\ge32$, transformer $d\ge64$), correlationally and partly circularly (runs stop 500 steps after generalisation). Circuit efficiency \citep{varma2023explaining} is only touched by E12/E12b: the logit spectrum is compact soon after the fit and the embedding's later (not at all for the wide MLP), which is compatible with a late-forming sparse circuit in the transformer and the narrow MLP, from five seeds and without intervention. Lazy-to-rich dynamics \citep{kumar2024grokking} is probed only through the MLP's initial scale (Sec.~4.5), where no effect was detected; with inconclusive slopes this is absence of evidence, not a refutation.

\textbf{Scope and failure cases.} The results are for one-hidden-layer MLPs and one-block transformers without LayerNorm (plus a depth-2 control with 3 seeds), widths $4$--$256$ (parameters below $10^6$), $p=59$ (mostly) and $97$, three operations, full-batch AdamW, and the stated step caps. Censoring matters: an MLP that does not generalise within 20k steps at $f=0.3$ generalises at 30--100k, and $f=0.2$ is untested beyond 100k. The 0.99 threshold is sensitive to plateaus at about 98.6\% validation accuracy, so only verdicts of the same class at 0.99 and 0.95 are counted. The transformer has large seed variance (standard deviation of $\log_{10}t_{\mathrm{gen}}$ 0.18--0.44 per cell), its three-seed cells cannot detect slopes below 0.3--0.7, and its learning rate was not tuned at weight decay 3. In E7 the learning-rate optimum sits at the top of the grid and $t_{\mathrm{gen}}$ at the logging floor; E14 removes both limits for the MLP at $f=0.3$--$0.4$, one weight decay. The $\pm0.15$ margin is a convention.

\textbf{Not done.} Mechanistic progress measures beyond E12 (a larger width grid, interventions), a transformer initialisation-scale sweep, a parametrisation sweep, $p=113$, transformers with other operations, widths above 256, depth above 2 and a cross-machine reproduction were not run. Split and initialisation noise are not separated (one split per seed).

\subsection*{AI use statement}
This study was carried out end to end by generative-AI agents (Claude models) organised in separate roles: one read the task and the related work and wrote a brief; one proposed the hypothesis, the predictions and the study plan and revised it after pilot results; one judged feasibility; one wrote the training and analysis code, ran the experiments and the statistics, and wrote the verdicts; one wrote the figure and table generators; one chose the venue format; one wrote this paper; and further agents reviewed the work independently, including an external-style review of a first version whose points led to the follow-up experiments E7--E13. The task statement was the only human input. Hypotheses, experiment design, implementation, interpretation of results, literature search (limited to the supplied starter list and look-ups by identifier), figures and writing were therefore all AI-generated; the authors of record take responsibility for the content, and every number in the paper is taken from files produced by the released code. No human annotation or human-subject data were used.

\subsection*{Reproducibility statement}
All data are generated in memory. The repository contains the code (\texttt{code/}), the per-run results (\texttt{results/}), the figure and table scripts (\texttt{figures/}, \texttt{tables/}) and \texttt{REPRODUCE.md}, which lists the environment and versions (Python 3.11.2, PyTorch 2.14.1 CPU), the seeding, the exact commands, the run counts (1\,852 runs in all) and the cost (an estimated 60.3 dedicated single-core CPU-hours, 198 contended wall-hours on a shared host, 0.1 s to 65 min per run), and the script behind each figure and table; \texttt{code/reproduce.sh} has entry points for a 10-run quick check, the tables and the figures. On the host used, every configuration computed twice reproduced $t_{\mathrm{fit}}$ and $t_{\mathrm{gen}}$ exactly (10 quick runs, 20 E6, 50 E9 and 9 E11 runs); CPU kernels are not guaranteed bit-identical across machines, which we did not test. No forbidden source was consulted. The related work is the supplied starter list plus five earlier papers (\citealp{liu2022towards}, \citealp{gromov2023grokking}, \citealp{thilak2022slingshot}, \citealp{zhong2023clock}, \citealp{merrill2023tale}), looked up by arXiv identifier after an external reviewer named them; only their abstracts were read, and no topical literature search was run. ICLR 2027 has no separate paper checklist; the AI-use statement above is the required disclosure.

\subsection*{Ethics statement}
All data are synthetic modular-arithmetic examples generated in memory; no human subjects are involved.

\bibliography{references}
\bibliographystyle{iclr2027_conference}

\appendix
\section{Appendix}
\subsection{Plan history and deviations}\label{app:plan}
The plan is delivered as \texttt{plan/PLAN.md} with its version history (\texttt{plan/HISTORY.md}). Predictions P1 and P2--P6 with their falsifiers were committed at 13:26 UTC (version 1), before the first training code (13:30) and any result, and were sharpened, not reversed, up to version 4.1 (13:51). The first E1 and E2-MLP sweeps were already running at about 14:10, but per-run launch times were not logged, so the ordering rests on the plan history and the deviations log. P1b and the separate treatment of $t_{\mathrm{fit}}$, $t_{\mathrm{gen}}$ and $\Delta$ in P2 were added in version 5 (14:44), after a two-seed pilot and after the main sweeps had started, so they are not pre-stated in the strict sense; P7 (version 6) is post hoc. The deviations (D1--D5 in \texttt{code/SPEC.md}) are: the width sweeps at $d\ge16$ (E1, E2-MLP) were launched before the specification check (D1), though all logs are stored and nothing was rerun; the 8-seed transformer analysis was pre-specified, while the 16-seed analysis was added after the seed standard deviation was seen (D4); and the sustained-crossing definition and the 0.95 co-primary threshold were adopted after validation-accuracy plateaus at about 98.6\% were seen (D3). Pre-stated are the width trend of the delay and of $t_{\mathrm{fit}}$ (P2), the capacity gate, the lr/wd exponents and the tuned-versus-untuned test with fresh seeds. Exploratory or post hoc are $f_{\mathrm{gen}}(w)$ and the separate $t_{\mathrm{gen}}$ trend, P7 and its held-out tests, the per-fraction slopes, the scope settings, E5-lite, E6, the decision rule and all follow-ups E7--E13 (requested by the review).

\subsection{Additional figures}

\begin{figure}[t]
\centering
\includegraphics[width=\linewidth]{f4_controls.pdf}
\caption{Controls (MLP unless stated). (a, b) Generalisation time against lr$\cdot$wd before and after multiplying by lr$\cdot$wd; each marker is a run (3 seeds per cell). (c, d) Baseline (lr $10^{-3}$, wd 1) against the per-width tuned setting on fresh seeds 3--5, disjoint from the seeds 0--2 used to select it; tuned log--log slope with 95\% CI; the tuned cell of each width is labelled, so the change from wd 1 at $d=16$ to wd 3 at $d\ge32$ is visible, and panel (c) also gives the slope within wd 3 alone ($d=32$--$128$). Transformer: three widths, $d=16$ tuned equals baseline.}
\label{fig:controls}
\end{figure}

\begin{figure}[t]
\centering
\includegraphics[width=\linewidth]{f5_scope.pdf}
\caption{Width dependence of fit time (a) and generalisation time at 0.99 (b) for other operations and prime (MLP; transformer dashed), medians with 95\% bootstrap CIs over seeds: addition at $p=59$ and $p=97$ 8 seeds, subtraction and $a^2+b^2$ 5 seeds, transformer at $p=97$ 3 seeds (no interval). Slopes are the log--log Tobit slopes from the analysis summaries.}
\label{fig:scope}
\end{figure}

\begin{figure}[t]
\centering
\includegraphics[width=\linewidth]{f1_curves.pdf}
\caption{Train and validation accuracy (upper rows) and total parameter norm (lower rows) against step for every seed of the width sweep ($f=0.5$, addition mod 59), at $d=16,64,256$ for the MLP (8 seeds) and the transformer (16 seeds; 8 at $d=256$). Open and filled markers: $t_{\mathrm{fit}}$ and $t_{\mathrm{gen}}$ (0.99).}
\label{fig:curves}
\end{figure}

\begin{figure}[t]
\centering
\includegraphics[width=\linewidth]{f3_outcome_map.pdf}
\caption{Outcome of each run by width and train fraction (share of seeds per class; numbers are seeds per cell; MLP and transformer, cap 20k and 25k steps; the long-cap grid is in \pfig{fig:fgen}). Runs that do not generalise within the cap are kept as ``fit, no generalisation in budget''; cells with 3 seeds give fractions, not intervals. Black bars in (a, b): $f_{\mathrm{gen}}(d)$, the smallest $f$ at which at least 2/3 of the seeds generalise (0.99). (c, d): Tobit log--log slope of $t_{\mathrm{gen}}$ (orange) and $t_{\mathrm{fit}}$ (grey) on width at each $f$ with 95\% CI; the shaded band is $\pm0.15$; the labels give generalising runs over runs.}
\label{fig:map}
\end{figure}

\begin{figure}[t]
\centering
\includegraphics[width=\linewidth]{f6_capacity_gate.pdf}
\caption{Share of seeds that fit against width: random labels ($f=0.3,0.5,0.7$) and true labels.}
\label{fig:capacity}
\end{figure}

\begin{figure}[t]
\centering
\includegraphics[width=\linewidth]{f7_gen_time_vs_f.pdf}
\caption{Generalisation time against width, one line per train fraction; cells with 3 seeds carry no interval.}
\label{fig:vsf}
\end{figure}

\subsection{Additional tables}
\input{../tables/t3b_verdicts_detail.tex}
\input{../tables/t1_inventory.tex}
\input{../tables/t2_e1_mlp.tex}
\input{../tables/t2_e1_tfm.tex}
\input{../tables/t4_slopes_by_f.tex}
\input{../tables/t5_timescale.tex}
\input{../tables/t6_capacity.tex}
\input{../tables/t7_key_numbers.tex}
\input{../tables/t8_threshold_rule.tex}
\input{../tables/t8b_threshold_rule_followups.tex}
\input{../tables/t9_fine_tfit.tex}
\input{../tables/t10_fourier.tex}
\input{../tables/t11_depth.tex}

\end{document}
